\documentclass[a4paper]{article} %

\date{9.07.2026} %

% Please, do not change any of the above lines

\usepackage{amssymb} % add other LaTeX packages you need here.

\begin{document}
\setcounter{page}{1} % Please, do not delete this line

% Your title
\title{Clinker equivalences}

% Authors
\author{S. Radeleczki, L. Veres, \\ %
       Miskolc University, Hungary\\ \\ % Affiliation 1
       J. J\"{a}rvinen \\ % If any other author with different Affilation
       LUT School of Engineering Science, Finland\\ \\ % Affiliation 2 (if needed)
       % New author \\
       % New affiliation \\
       % Add authors and affiliation as needed
       \texttt{sandor.radeleczki@uni-miskolc.hu} % Only one corresponding e-mail, please
       }%

\maketitle



We study the maximal relations $\varrho\subseteq U\times U$ that define atomic
Boolean concept lattices $\mathfrak{V}(U,U,\varrho)$. The complement
$R=\left(  U\times U\right)  \setminus\varrho$ of such a relation $\varrho$
can be characterized as a reflexive relation with the property that the sets
$R^{-1}(x):=$ $\{y\in U\mid(y,x)\in R\}$ form an irredundant covering of $U$.
It was proved by J. J\"{a}rvinen and S. Radeleczki that the rough sets induced
by a relation $R$ with this property form an algebraic regular double Stone
lattice, i.e. they have the same structure as in the case of an equivalence
relation. Therefore, such a reflexive relation $R$ can be viewed as a new
generalization of the notion of an equivalence relation, however, in the case
when $R$ is symmetric or transitive, it becomes an ordinary equivalence. The
term "clinker equivalence" was chosen as its name, because the geometric form
of the covering $U$ $=\{R^{-1}(x)\mid x\in U\}$ resembles the sail of Viking
ship, specifically referencing the historical clinker-built method of Norse
shipbuilding. Here, we present some examples and applications of these
relations. We show that they can be interpreted as special directed similarity
relations, and can serve as useful tools in classification problems and
concept approximations.

\begin{thebibliography}{99}                                                                                               %
\bibitem[1]{}Stephen D. Comer, \textit{ On connections between information
systems, rough sets, and algebraic logic}, Algebraic methods in Logic and
Computer Science, 1993, pp. 117-124.

\bibitem[2]{}B. Ganter and R. Wille, \textit{ Formal concept analysis:
Mathematical foundations}, Springer, Berlin-Heidelberg, 1999.

\bibitem[3]{}J. J\"{a}rvinen, S. Radeleczki and L. Veres,\textit{ Rough sets
determined by quasiorders}, Order, \textbf{26}, (2009) 337-355.

\bibitem[4]{}J. J\"{a}rvinen and S. Radeleczki, \textit{ Representation of
Nelson algebras by rough sets determined by quasiorder}, Algebra Universalis
\textbf{66} (2011), 163-179.

\bibitem[5]{}J. J\"{a}rvinen and S. Radeleczki, \textit{ Rough sets determined
by tolerances}, Internat. J. Approx. Reason. \textbf{55} (2014), 1419-1438.

\bibitem[6]{}J. J\"{a}rvinen and S. Radeleczki, \textit{Kleene and Stone
algebras induced by reflexive relations}, Internat. J. Approx. Reason.
accepoted for publication.

\bibitem[7]{}J. J\"{a}rvinen and S. Radeleczki, \textit{Clinker equivalences
and their applications}, Information Science, submitted (2026).

\bibitem[8]{}T. Katrin\'{a}k, \textit{ The structure of didtributive double
p-algebras. Regularity and congruences} Algebra Universalis \textbf{3}, (1973) 238-246.

\bibitem[9]{}Z. Pawlak,\textit{ Rough sets}, International J. Comput. Inform.
Sci. \textbf{11}, (1982) 341-356.

\bibitem[10]{}J. Pomykala and J. A. Pomykala,\textit{ The Stone algebra of
rough sets}, Bull. Polish Acad. Sci. Math. \textbf{36},(1988) 495-512.

\bibitem[11]{}H.P. Sankapanavar, \textit{ Pseudocomplemented Ockham and De
Morgan algebras}, MLQ Math. Log. Q. \textbf{32}, (1986), 385-394.

\bibitem[12]{}R. Slowinski, D. Vanderpooten, \textit{Similarity relations as a
basis for rough approximations}, Adv Mach Intell Soft Comput \textbf{4} (1997) 17-33.

\bibitem[13]{}T. T\'{o}th, S. Radeleczki, L. Veres and A. K\"{o}rei, \textit{A
new mathematical approach to supporting group technology}, European J.
Industrial Engineering, Vol. 8, No. \textbf{5,} (2014) 716-737.

\bibitem[14]{}J. Varlet, \textit{ A regular variety of type (2,2,1,1,0,0)},
Algebra Universalis \textbf{2}, (1972) 218-223.
\end{thebibliography}


\end{document}
